\subsubsection*{命题 6.3.}
存在同构

\begin{equation}\tag{6.3.1}
W_S(1)\simeq
\operatorname{Hom}_{\mathbb F_q}\!\left(
H^{n+1}(\mathbb P^{n+1},\mathcal O_{\mathbb P}(-d)),
\mathbb F_q\right)
\end{equation}

在此同构下, 算子

\begin{equation}\tag{6.3.2}
\psi_q\circ(X_0f(X))^{q-1}:W_S(1)\longrightarrow W_S(1)
\end{equation}

是下列算子的转置

\begin{equation}\tag{6.3.3}
f(X)^{q-1}\mathfrak F_q:
H^{n+1}(\mathbb P^{n+1},\mathcal O_{\mathbb P}(-d))
\longrightarrow
H^{n+1}(\mathbb P^{n+1},\mathcal O_{\mathbb P}(-d)).
\end{equation}

\noindent\emph{证明.}
利用 $\mathbb P^{n+1}$ 的标准覆盖计算
$H^{n+1}(\mathbb P^{n+1},\mathcal O_{\mathbb P}(-d))$, 得到

\begin{equation}\tag{6.3.4}
\begin{aligned}
H^{n+1}(\mathbb P^{n+1},\mathcal O_{\mathbb P}(-d))
&=\text{以下列元素为基的 $\mathbb F_q$ 向量空间的商}\\
&\quad X_1^{W_1}\cdots X_{n+2}^{W_{n+2}},
\qquad W_i\in\mathbb Z,\quad\sum_iW_i=-d,\\
&\quad\text{模去以下列单项式为基的 $\mathbb F_q$ 子空间}\\
&\quad X_1^{W_1}\cdots X_{n+2}^{W_{n+2}}
\text{ 满足 }\exists i\text{ 使 }W_i\geq0.
\end{aligned}
\end{equation}

另一方面, $W_S(1)=R_S(1)\otimes_{\mathcal O}\mathbb F_q$, 且

\begin{equation}\tag{6.3.5}
\begin{aligned}
R_S(1)
&=\text{以下列单项式为基的自由 $\mathcal O$ 模}\\
&\quad
(\pi X_0)X_1^{U_1}\cdots X_{n+2}^{U_{n+2}},
\qquad
\sum_iU_i=d,\quad U_i>0\ \forall i.
\end{aligned}
\end{equation}

配对

\begin{equation}\tag{6.3.6}
(\ ,\ ):
W_S(1)\times
H^{n+1}(\mathbb P^{n+1},\mathcal O_{\mathbb P}(-d))
\longrightarrow\mathbb F_q
\end{equation}

定义为

\begin{equation}\tag{6.3.7}
\left(
\operatorname{class}\!\left(
\pi X_0X_1^{U_1}\cdots X_{n+2}^{U_{n+2}}
\right),
X_1^{W_1}\cdots X_{n+2}^{W_{n+2}}
\right)
=
\begin{cases}
1,&U_i=-W_i\text{ 对每个 }i,\\
0,&\text{其他情形},
\end{cases}
\end{equation}

显然是 $\mathbb F_q$ 向量空间的完美对偶. 立即可验证
$\psi_q\circ(X_0f(X))^{q-1}$ 是 $f(X)^{q-1}\mathfrak F_q$ 的转置,
这就证明了命题~6.3.

\medskip
\noindent\emph{底本注.}
在 (6.3.4) 中, 印本写作 $H^{i+1}$, 而命题~6.3、(6.3.3)
以及引出该计算的句子均要求为 $H^{n+1}$.
校订法文层和英文层采用 $n+1$; 照录法文层保留印本指数.

\section*{参考文献}

\begin{enumerate}[label={[\arabic*]},leftmargin=2.4em]
\item J.~Ax, \emph{Zeroes of Polynomials Over Finite Fields},
  \emph{Amer. J. Math.} \textbf{86} (1964), 255--261.
\item B.~Dwork, \emph{On the Rationality of the Zeta Function of an Algebraic
  Variety}, \emph{Amer. J. Math.} \textbf{82} (1960), 631--648.
\item B.~Dwork, \emph{An Algebraic Variety}, 载于 \emph{Proceedings of the
  Woods Hole Summer Institute 1964} (油印本).
\item B.~Dwork, \emph{On the Zeta Function of a Hypersurface II},
  \emph{Ann. of Math.} (2) \textbf{80} (1964), 227--299.
\item A.~Grothendieck, \emph{Le Groupe de Brauer III}, 载于
  \emph{Dix Exposés sur la Cohomologie des Schémas},
  North-Holland Publishing Co., 1968.
\item J.-P.~Serre, \emph{Groupes algébriques et corps de classes},
  Hermann, Paris, 1959.
\item J.-P.~Serre, \emph{Endomorphismes complètement continus des espaces de
  Banach $p$-adiques}, \emph{Publ. Math. I.H.E.S.} no.~12 (1962).
\item C.~Chevalley, \emph{Démonstration d'une hypothèse de M.~Artin},
  \emph{Abh. Math. Sem. Univ. Hamburg} \textbf{11} (1936), 73--75.
\end{enumerate}
